% ============================================================
%  GUÍA 11: Teoremas, Monotonía y Concavidad
%  Curso de Introducción a la Universidad — Precálculo y Cálculo I
%  Autor: Ing. Gabriel Astudillo
%  Clase cubierta: 11
% ============================================================

\documentclass[12pt, a4paper]{article}

% ── Paquetes ─────────────────────────────────────────────────
\usepackage{mathwizards-guia}
\mwguia{Guía 11 — Teoremas, Monotonía y Concavidad}{Ing. Gabriel Astudillo $\cdot$ Curso Preuniversitario}

\begin{document}

% ── Portada / Título ─────────────────────────────────────────
\mwportada{Guía de Estudio 11}{Teoremas, Monotonía y Concavidad}{Teoremas de Rolle y del Valor Medio, Criterios de la Primera y Segunda Derivada, Trazado de Curvas}

\vspace{4pt}
\begin{cuadroinfo}
  \small
  \textbf{Presentación:} Esta guía aplica la derivada al análisis de funciones: los teoremas de Rolle y del Valor Medio, los criterios de la primera y segunda derivada, los puntos críticos y de inflexión y el trazado completo de curvas. Incluye teoría, ejercicios resueltos paso a paso y una amplia batería de propuestos con su clave de respuestas.
  \\[4pt]
  \textbf{Nivel de Dificultad:}\quad
  \facil{} Básico / Directo \quad
  \medio{} Intermedio \quad
  \dificil{} Desafiante \quad
  \muyDificil{} Muy Difícil / Reto
\end{cuadroinfo}

\vspace{8pt}

\section{Máximos y Mínimos Locales}
% ============================================================

\subsection{Conceptos fundamentales}

\begin{cuadroteoria}
  \textbf{Valor máximo local:} $f(c)$ es máximo local si $f(c) \geq f(x)$ para todo $x$ cercano a $c$.

  \textbf{Valor mínimo local:} $f(c)$ es mínimo local si $f(c) \leq f(x)$ para todo $x$ cercano a $c$.

  \textbf{Punto crítico:} Un punto $c$ del dominio de $f$ donde $f'(c) = 0$ o $f'(c)$ no existe.

  \textbf{Teorema (condición necesaria):} Si $f$ tiene un extremo local en $c$ y $f'(c)$ existe, entonces $f'(c) = 0$. Los extremos locales ocurren en puntos críticos.
\end{cuadroteoria}

\begin{cuadrosolucion}
  \textbf{Prueba de la primera derivada:}
  \begin{itemize}[leftmargin=*]
    \item Si $f'$ cambia de $+$ a $-$ al pasar por $c$: máximo local en $c$.
    \item Si $f'$ cambia de $-$ a $+$ al pasar por $c$: mínimo local en $c$.
    \item Si $f'$ no cambia de signo: ni máximo ni mínimo.
  \end{itemize}
\end{cuadrosolucion}
\newpage
\begin{cuadroextra}
  \textbf{Prueba de la segunda derivada:}
  Sea $c$ un punto crítico donde $f'(c) = 0$:
  \begin{itemize}[leftmargin=*]
    \item Si $f''(c) < 0$: $f$ es cóncava hacia abajo en $c \;\Rightarrow\; f(c)$ es un \textbf{máximo local}.
    \item Si $f''(c) > 0$: $f$ es cóncava hacia arriba en $c \;\Rightarrow\; f(c)$ es un \textbf{mínimo local}.
    \item Si $f''(c) = 0$: la prueba \textbf{no es concluyente} (usar la primera derivada).
  \end{itemize}
\end{cuadroextra}

\begin{cuadroadvertencia}
  \textbf{Teorema de Fermat:} Si $f$ es continua en $[a, b]$, entonces $f$ alcanza un valor máximo y un valor mínimo absolutos en ese intervalo. Para encontrarlos:
  \begin{enumerate}
    \item Halla los puntos críticos en $(a, b)$.
    \item Evalúa $f$ en los puntos críticos y en los extremos $a$ y $b$.
    \item El mayor valor es el máximo absoluto; el menor, el mínimo absoluto.
  \end{enumerate}
\end{cuadroadvertencia}

\subsection{Ejercicios resueltos}

\textbf{Ejemplo R1.} Encuentra los puntos críticos de $f(x) = x^3 - 3x^2 + 2$.

\begin{cuadrosolucion}
  \textbf{Solución:}
  \begin{enumerate}
    \item $f'(x) = 3x^2 - 6x = 3x(x - 2)$.
    \item $f'(x) = 0 \Rightarrow 3x(x - 2) = 0 \Rightarrow x = 0$ o $x = 2$.
  \end{enumerate}
  Puntos críticos: $x = 0$ y $x = 2$.
\end{cuadrosolucion}

\textbf{Ejemplo R2.} Clasifica los puntos críticos de $f(x) = x^3 - 3x^2$ usando la prueba de la primera derivada.

\begin{cuadrosolucion}
  \textbf{Solución:} $f'(x) = 3x(x - 2)$. Los puntos críticos son $x = 0$ y $x = 2$. Analizamos signos:
  \begin{center}
    \small
    \begin{tabular}{lccc}
      \toprule
      \textbf{Intervalo} & $(-\infty, 0)$ & $(0, 2)$    & $(2, \infty)$ \\
      \midrule
      Signo de $f'$      & $+$            & $-$         & $+$           \\
      Comportamiento     & creciente      & decreciente & creciente     \\
      \bottomrule
    \end{tabular}
  \end{center}
  En $x = 0$: $f'$ cambia de $+$ a $-$, hay máximo local $f(0) = 0$.

  En $x = 2$: $f'$ cambia de $-$ a $+$, hay mínimo local $f(2) = -4$.
\end{cuadrosolucion}
\newpage
\textbf{Ejemplo R3.} Encuentra los extremos absolutos de $f(x) = x^2 - 4x + 5$ en $[0, 5]$.

\begin{cuadrosolucion}
  \textbf{Solución:}
  \begin{enumerate}
    \item Derivada: $f'(x) = 2x - 4$. Cero: $x = 2$ (en el intervalo).
    \item Evaluamos:
          \begin{itemize}
            \item $f(0) = 5$
            \item $f(2) = 4 - 8 + 5 = 1$
            \item $f(5) = 25 - 20 + 5 = 10$
          \end{itemize}
    \item Mínimo absoluto: $1$ (en $x = 2$). Máximo absoluto: $10$ (en $x = 5$).
  \end{enumerate}
\end{cuadrosolucion}

\textbf{Ejemplo R4.} Encuentra los puntos críticos de $f(x) = \sqrt[3]{(x - 1)^2}$.

\begin{cuadrosolucion}
  \textbf{Solución:} Reescribimos $f(x) = |x - 1|^{2/3}$. La derivada:
  $$f'(x) = \frac{2}{3}(x - 1)^{-1/3} = \frac{2}{3\sqrt[3]{x - 1}}.$$
  La derivada nunca es cero, pero no existe en $x = 1$. Por tanto, $x = 1$ es punto crítico (donde la función tiene un pico).
\end{cuadrosolucion}

\subsection{Ejercicios propuestos}

\begin{enumerate}[series=ejercicios, leftmargin=*]
  \item \facil{} Encuentra los puntos críticos de $f(x) = x^2 - 6x + 10$.

  \item \facil{} Encuentra los puntos críticos de $f(x) = x^4 - 4x^2$.

  \item \facil{} Encuentra los extremos absolutos de $f(x) = x^2$ en $[-2, 1]$.

  \item \facil{} Encuentra los extremos absolutos de $f(x) = x + 1$ en $[0, 4]$.

  \item \medio{} Clasifica los puntos críticos de $f(x) = 2x^3 - 3x^2 - 36x$.

  \item \medio{} Encuentra los extremos absolutos de $f(x) = x^3 - 3x + 1$ en $[0, 3]$.

  \item \medio{} Encuentra los puntos críticos de $f(x) = \dfrac{x}{x^2 + 1}$.

  \item \medio{} Determina los máximos y mínimos locales de $f(x) = \sin x$ en $[0, 2\pi]$.

  \item \dificil{} Encuentra los valores de $a$ para que $f(x) = ax^3 - 6x^2$ tenga un punto crítico en $x = 2$.

  \item \dificil{} Determina los extremos absolutos de $f(x) = x\sqrt{9 - x^2}$ en $[-3, 3]$.

  \item \dificil{} Encuentra los extremos absolutos de $f(x) = \dfrac{x - 1}{x + 1}$ en $[0, 4]$.

  \item \dificil{} Prueba que $f(x) = x^3 + x + 1$ tiene exactamente un punto crítico (y que es un mínimo local). ¿Cuál es?
\end{enumerate}

\newpage

% ============================================================

\section{Teorema de Rolle y Teorema del Valor Medio}
% ============================================================

\subsection{Enunciados}

\begin{cuadroteoria}
  \textbf{Teorema de Rolle:} Si $f$ es continua en $[a, b]$, derivable en $(a, b)$, y $f(a) = f(b)$, entonces existe al menos un $c \in (a, b)$ tal que $f'(c) = 0$.

  \textbf{Teorema del Valor Medio (TVM):} Si $f$ es continua en $[a, b]$ y derivable en $(a, b)$, entonces existe al menos un $c \in (a, b)$ tal que:
  $$f'(c) = \frac{f(b) - f(a)}{b - a}$$
\end{cuadroteoria}

\begin{cuadrosolucion}
  \textbf{Interpretación geométrica del TVM:} existe un punto donde la tangente es paralela a la recta secante que une los extremos del intervalo.
\end{cuadrosolucion}

\subsection{Ejercicios resueltos}

\textbf{Ejemplo R1.} Verifica el Teorema de Rolle para $f(x) = x^2 - 4x$ en $[0, 4]$ y encuentra $c$.

\begin{cuadrosolucion}
  \textbf{Solución:}
  \begin{itemize}
    \item $f$ es polinomial: continua y derivable en todo $\mathbb{R}$.
    \item $f(0) = 0$; $f(4) = 16 - 16 = 0$. Se cumple $f(0) = f(4)$.
    \item $f'(x) = 2x - 4$. Igualamos a cero: $2x - 4 = 0 \Rightarrow x = 2$.
  \end{itemize}
  $c = 2$ pertenece a $(0, 4)$ y cumple $f'(c) = 0$.
\end{cuadrosolucion}

\textbf{Ejemplo R2.} Encuentra $c$ según el TVM para $f(x) = x^3$ en $[1, 2]$.

\begin{cuadrosolucion}
  \textbf{Solución:}
  \begin{itemize}
    \item Pendiente secante: $\dfrac{f(2) - f(1)}{2 - 1} = \dfrac{8 - 1}{1} = 7$.
    \item Derivada: $f'(x) = 3x^2$. Igualamos: $3c^2 = 7 \Rightarrow c^2 = \frac{7}{3} \Rightarrow c = \sqrt{\frac{7}{3}} \approx 1{,}53 \in (1, 2)$.
  \end{itemize}
\end{cuadrosolucion}

\textbf{Ejemplo R3.} Un automóvil acelera de $0$ a $100$ km/h en $10$ segundos. ¿Qué te dice el TVM?

\begin{cuadrosolucion}
  \textbf{Solución:} La velocidad promedio es $\dfrac{100 - 0}{10}$ km/h/s $= 10$ km/h/s. El TVM garantiza que en algún instante la velocidad instantánea (derivada de la posición) fue exactamente esa velocidad promedio. (Y de hecho, si la función de velocidad es continua, también se puede aplicar a la velocidad.)
\end{cuadrosolucion}

\newpage

\subsection{Ejercicios propuestos}

\begin{enumerate}[resume=ejercicios, leftmargin=*]
  \item \facil{} Verifica el Teorema de Rolle para $f(x) = x^2 - 1$ en $[-1, 1]$ y encuentra $c$.

  \item \facil{} Verifica el Teorema de Rolle para $f(x) = \sin x$ en $[0, \pi]$ y encuentra $c$.

  \item \facil{} Encuentra $c$ según el TVM para $f(x) = 2x + 1$ en $[0, 3]$.

  \item \facil{} Encuentra $c$ según el TVM para $f(x) = x^2$ en $[2, 4]$.

  \item \medio{} Verifica si se puede aplicar Rolle a $f(x) = x^2 - 4x + 3$ en $[1, 3]$ y halla $c$.

  \item \medio{} Encuentra $c$ según el TVM para $f(x) = \sqrt{x}$ en $[1, 4]$.

  \item \medio{} Un objeto se mueve con posición $s(t) = t^2 + 3t$ metros entre $t = 0$ y $t = 2$ s. ¿Cuál es la velocidad promedio? ¿En qué instante $t$ la velocidad instantánea coincide con la promedio?

  \item \medio{} Demuestra usando el TVM que si $f'(x) = 0$ para todo $x$, entonces $f$ es constante en el intervalo.

  \item \dificil{} Aplica el TVM a $f(x) = \ln x$ en $[1, e]$ y encuentra $c$.

  \item \dificil{} Sea $f$ continua y derivable con $f(1) = 2$ y $f(5) = 10$. ¿Qué puedes afirmar sobre $f'(c)$ para algún $c \in (1, 5)$?

  \item \dificil{} Prueba que la ecuación $x^3 - 3x + 1 = 0$ tiene al menos una raíz en $(0, 1)$ (pista: usa TVM sobre alguna función conveniente, o el TVI).

  \item \dificil{} Explica por qué el TVM no se puede aplicar a $f(x) = |x|$ en $[-1, 1]$.
\end{enumerate}

%\newpage

% ============================================================

\section{Análisis Completo de una Gráfica}
% ============================================================

\subsection{Concavidad y puntos de inflexión}

\begin{cuadroteoria}
  \textbf{Concavidad:}
  \begin{itemize}[leftmargin=*]
    \item Si $f''(x) > 0$ en un intervalo: cóncava hacia arriba (la curva se abre como una taza).
    \item Si $f''(x) < 0$ en un intervalo: cóncava hacia abajo (la curva se abre como una montaña).
  \end{itemize}

  \textbf{Punto de inflexión:} punto donde la concavidad cambia (donde $f''(x) = 0$ o no existe, y cambia de signo).
\end{cuadroteoria}

\begin{cuadropasos}
  \textbf{Pasos para graficar una función con análisis:}
  \begin{enumerate}
    \item Determina el \textbf{dominio}, y encuentra las \textbf{intersecciones} con los ejes.
    \item Analiza \textbf{simetrías} (par/impar).
    \item Encuentra las \textbf{asíntotas} (verticales, horizontales, oblicuas).
    \item Calcula $f'$ y determina \textbf{intervalos de crecimiento/decrecimiento} y extremos.
    \item Calcula $f''$ y determina \textbf{concavidad} y puntos de inflexión.
    \item Esboza la gráfica con toda la información.
  \end{enumerate}
\end{cuadropasos}

\subsection{Ejercicios resueltos}

\textbf{Ejemplo R1.} Determina intervalos de crecimiento, concavidad y puntos de inflexión de $f(x) = x^3 - 6x^2 + 9x$.

\begin{cuadrosolucion}
  \textbf{Solución:}
  \begin{enumerate}
    \item Dominio: $\mathbb{R}$.
    \item $f'(x) = 3x^2 - 12x + 9 = 3(x^2 - 4x + 3) = 3(x - 1)(x - 3)$.
          \begin{itemize}
            \item $f'(x) > 0$ en $(-\infty, 1) \cup (3, \infty)$: creciente.
            \item $f'(x) < 0$ en $(1, 3)$: decreciente.
            \item $x = 1$: máximo local ($f(1) = 4$). $x = 3$: mínimo local ($f(3) = 0$).
          \end{itemize}
    \item $f''(x) = 6x - 12 = 6(x - 2)$.
          \begin{itemize}
            \item $f''(x) < 0$ en $(-\infty, 2)$: cóncava hacia abajo.
            \item $f''(x) > 0$ en $(2, \infty)$: cóncava hacia arriba.
            \item Punto de inflexión en $x = 2$: $f(2) = 8 - 24 + 18 = 2$. Punto $(2, 2)$.
          \end{itemize}
  \end{enumerate}
\end{cuadrosolucion}

\textbf{Ejemplo R2.} Analiza y grafica $f(x) = \dfrac{x^2}{x - 1}$.

\begin{cuadrosolucion}
  \textbf{Solución:}
  \begin{enumerate}
    \item Dominio: $x \neq 1$.
    \item Interceptos: $x = 0$, $y = 0$. Punto $(0, 0)$.
    \item Asíntotas: A.V. en $x = 1$ (el denominador se anula y el numerador no). A.O.: dividiendo $x^2$ entre $x - 1$: $x + 1 + \frac{1}{x-1}$, así que $y = x + 1$ es asíntota oblicua.
    \item $f'(x) = \dfrac{2x(x - 1) - x^2}{(x - 1)^2} = \dfrac{x^2 - 2x}{(x - 1)^2} = \dfrac{x(x - 2)}{(x - 1)^2}$.
          \begin{itemize}
            \item Signo: positiva en $(-\infty, 0) \cup (2, \infty)$; negativa en $(0, 1) \cup (1, 2)$.
            \item $x = 0$: máximo local $f(0) = 0$. $x = 2$: mínimo local $f(2) = 4$.
          \end{itemize}
    \item $f''(x) = \dfrac{2}{(x - 1)^3}$ (verificar): es negativa en $(-\infty, 1)$ y positiva en $(1, \infty)$. Concavidad hacia abajo a la izquierda de 1, hacia arriba a la derecha.
  \end{enumerate}
\end{cuadrosolucion}

\newpage

\subsection{Ejercicios propuestos}

\begin{enumerate}[resume=ejercicios, leftmargin=*]
  \item \facil{} Determina la concavidad de $f(x) = x^2$ en todo su dominio.

  \item \facil{} Determina la concavidad de $f(x) = -x^2 + 3x$ en todo su dominio.

  \item \facil{} Encuentra los puntos de inflexión de $f(x) = x^3$.

  \item \facil{} Determina los intervalos de crecimiento de $f(x) = x^2 - 4x$.

  \item \medio{} Para $f(x) = x^3 - 3x$:
        \begin{enumerate}[label=\alph*)]
          \item Encuentra los puntos críticos y clasifícalos.
          \item Intervalos de crecimiento/decrecimiento.
          \item Intervalos de concavidad y puntos de inflexión.
        \end{enumerate}

  \item \medio{} Para $f(x) = \dfrac{1}{x}$:
        \begin{enumerate}[label=\alph*)]
          \item Determina los intervalos de crecimiento/decrecimiento.
          \item Determina la concavidad en cada intervalo del dominio.
          \item Identifica las asíntotas.
        \end{enumerate}

  \item \medio{} Encuentra los puntos de inflexión de $f(x) = x^4 - 6x^2$.

  \item \medio{} Analiza la gráfica de $f(x) = x e^{-x}$: dominio, extremos, crecimiento, concavidad y asíntota horizontal.

  \item \dificil{} Analiza y grafica $f(x) = \dfrac{2x}{x^2 - 1}$: dominio, interceptos, asíntotas, extremos, concavidad.

  \item \dificil{} Determina el valor de $a$ para que $f(x) = x^3 - ax^2 + 3$ tenga un punto de inflexión en $x = 1$.

  \item \dificil{} Analiza $f(x) = \ln(x^2 + 1)$: crecimiento, concavidad, extremos y puntos de inflexión.
\end{enumerate}

\newpage

% ============================================================

% ──────────────────────────────────────────────────────────
%  NUEVOS EJERCICIOS PROPUESTOS (15) — INSERCIÓN MANUAL
% ──────────────────────────────────────────────────────────
\subsection{Ejercicios propuestos: Rolle, TVM y Extremos}

\begin{enumerate}[resume=ejercicios, leftmargin=*]
  \item \facil{} Halla los puntos críticos de $f(x) = x^2 - 6x + 5$.

  \item \medio{} Verifica las hipótesis del teorema de Rolle para $f(x) = x^2 - 4x$ en $[0, 4]$ y halla el valor de $c$.

  \item \medio{} Aplica el Teorema del Valor Medio a $f(x) = x^2$ en $[1, 3]$ y halla el valor de $c$.

  \item \medio{} Determina los intervalos de crecimiento de $f(x) = x^3 - 3x$.

  \item \dificil{} Clasifica los extremos locales de $f(x) = x^3 - 3x^2 + 2$ usando el criterio de la primera derivada.

  \item \dificil{} Halla los puntos de inflexión de $f(x) = x^3 - 3x^2 + 1$.

  \item \dificil{} Determina los intervalos de concavidad y los puntos de inflexión de $f(x) = x^4 - 4x^3$.

  \item \muyDificil{} Aplica el Teorema del Valor Medio a $f(x) = \sqrt{x}$ en $[1, 4]$ y halla el valor de $c$.
\end{enumerate}

\subsection{Ejercicios propuestos: Análisis de Gráficas y Retos}

\begin{enumerate}[resume=ejercicios, leftmargin=*]
  \item \facil{} Halla el punto crítico de $f(x) = x^2 + 2x$.

  \item \medio{} Determina si $f(x) = x^3$ tiene extremos locales. Justifica.

  \item \medio{} Realiza el trazado completo de $f(x) = x^2 - 4x + 3$: dominio, cortes con los ejes, vértice y concavidad.

  \item \dificil{} Halla los valores extremos absolutos de $f(x) = x^2 - 4x + 3$ en el intervalo $[0, 5]$.

  \item \dificil{} Analiza $f(x) = \dfrac{x^2}{x - 1}$: asíntotas, puntos críticos y clasificación de extremos.

  \item \muyDificil{} Realiza el análisis completo de $f(x) = x^3 - 6x^2 + 9x + 1$: extremos, monotonía, concavidad y puntos de inflexión.

  \item \muyDificil{} \textbf{Optimización:} Maximiza el área de un rectángulo cuyo perímetro es 40 m.
\end{enumerate}
\newpage
% ============================================================
\ifmwclaves
\section{Clave de Respuestas}
% ============================================================

\subsection*{Sección 1: Máximos y Mínimos (Ejercicios 1--12)}

\begin{enumerate}[series=respuestas, leftmargin=*, label=\textbf{\arabic*.}]
  \item $f'(x) = 2x - 6$, punto crítico: $x = 3$.
  \item $f'(x) = 4x^3 - 8x = 4x(x^2 - 2)$. Puntos críticos: $x = 0$, $x = \pm\sqrt{2}$.
  \item Mínimo absoluto en $(0, 0)$: $0$. Máximo absoluto en $(-2, 4)$: $4$.
  \item Mínimo absoluto en $(0, 1)$: $1$. Máximo absoluto en $(4, 5)$: $5$.
  \item $f'(x) = 6x^2 - 6x - 36 = 6(x^2 - x - 6) = 6(x - 3)(x + 2)$. Ceros: $x = 3$ y $x = -2$. Tabla de signos: $f'$ positiva en $(-\infty, -2)$, negativa en $(-2, 3)$, positiva en $(3, \infty)$. En $x = -2$: máximo local. En $x = 3$: mínimo local.
  \item $f'(x) = 3x^2 - 3$, cero en $x = \pm 1$. En $[0, 3]$ solo $x = 1$ está adentro. $f(0) = 1$, $f(1) = -1$, $f(3) = 19$. Mínimo: $-1$ en $x = 1$; máximo: $19$ en $x = 3$.
  \item $f'(x) = \dfrac{1 - x^2}{(x^2 + 1)^2}$. Ceros: $x = \pm 1$.
  \item $f'(x) = \cos x$. Ceros en $x = \frac{\pi}{2}$ (máximo, $f = 1$) y $x = \frac{3\pi}{2}$ (mínimo, $f = -1$). Además, los extremos $0$ y $2\pi$ dan $f = 0$.
  \item $f'(x) = 3ax^2 - 12x$. En $x = 2$: $f'(2) = 3a(4) - 24 = 12a - 24 = 0 \Rightarrow a = 2$.
  \item $f'(x) = \sqrt{9 - x^2} + x \cdot \frac{-x}{\sqrt{9-x^2}} = \frac{9 - 2x^2}{\sqrt{9-x^2}}$. Ceros: $x = \pm \frac{3\sqrt{2}}{2}$. Evaluando en esos y en los extremos: máximo $\frac{9}{2}$ en $x = \pm \frac{3\sqrt{2}}{2}$; mínimo $0$ en los extremos $x = \pm 3$.
  \item $f'(x) = \dfrac{(x + 1) - (x - 1)}{(x + 1)^2} = \dfrac{2}{(x + 1)^2} > 0$ siempre, creciente. Mínimo en $x = 0$: $f(0) = -1$. Máximo en $x = 4$: $f(4) = \frac{3}{5}$.
  \item $f'(x) = 3x^2 + 1 > 0$ siempre, no hay puntos críticos (la derivada nunca se anula). $f$ es estrictamente creciente en todo $\mathbb{R}$.
\end{enumerate}


\subsection*{Sección 2: Rolle y TVM (Ejercicios 13--24)}

\begin{enumerate}[resume=respuestas, leftmargin=*, label=\textbf{\arabic*.}]
  \item $f(-1) = 0$, $f(1) = 0$. $f'(x) = 2x = 0 \Rightarrow c = 0$.
  \item $f(0) = 0$, $f(\pi) = 0$. $f'(x) = \cos x = 0 \Rightarrow c = \frac{\pi}{2}$.
  \item $f$ es lineal con pendiente 2, así que $c$ puede ser cualquier punto en $(0, 3)$ (por ejemplo $c = 1$).
  \item $\frac{f(4) - f(2)}{4 - 2} = \frac{16 - 4}{2} = 6$. $f'(x) = 2x = 6 \Rightarrow c = 3$.
  \item $f(1) = 1 - 4 + 3 = 0$, $f(3) = 9 - 12 + 3 = 0$. $f$ es continua y derivable. $f'(x) = 2x - 4 = 0 \Rightarrow c = 2 \in (1, 3)$.
  \item Pendiente secante: $\frac{2 - 1}{4 - 1} = \frac{1}{3}$. $f'(x) = \frac{1}{2\sqrt{x}} = \frac{1}{3} \Rightarrow 2\sqrt{c} = 3 \Rightarrow c = \frac{9}{4}$.
  \item Posición promedio: $\frac{s(2) - s(0)}{2} = \frac{10 - 0}{2} = 5$ m/s. $v(t) = s'(t) = 2t + 3 = 5 \Rightarrow t = 1$ s.
  \item Por el TVM, para cualquier $a < b$: $f'(c) = \frac{f(b) - f(a)}{b - a} = 0 \Rightarrow f(b) = f(a)$. Como $a$ y $b$ son arbitrarios, $f$ es constante.
  \item $\frac{f(e) - f(1)}{e - 1} = \frac{1 - 0}{e - 1} = \frac{1}{e - 1}$. $f'(x) = \frac{1}{x} = \frac{1}{e - 1} \Rightarrow c = e - 1 \approx 1{,}718 \in (1, e)$.
  \item TVM: existe $c \in (1, 5)$ con $f'(c) = \frac{10 - 2}{5 - 1} = 2$.
  \item Sea $f(x) = x^3 - 3x + 1$. $f(0) = 1$, $f(1) = -1$. Por TVI existe al menos una raíz en $(0,1)$. (También se puede argumentar con el TVM sobre una primitiva.)
  \item $f(x) = |x|$ no es derivable en $x = 0 \in (-1, 1)$, que es un punto interior del intervalo. No se cumplen las hipótesis del TVM.
\end{enumerate}


\subsection*{Sección 3: Análisis de Gráficas (Ejercicios 25--35)}

\begin{enumerate}[resume=respuestas, leftmargin=*, label=\textbf{\arabic*.}]
  \item $f''(x) = 2 > 0$: cóncava hacia arriba en todo $\mathbb{R}$.
  \item $f''(x) = -2 < 0$: cóncava hacia abajo en todo $\mathbb{R}$.
  \item $f''(x) = 6x = 0 \Rightarrow x = 0$; cambia de signo (negativa a la izquierda, positiva a la derecha). Punto de inflexión en $(0, 0)$.
  \item $f'(x) = 2x - 4$; creciente en $(2, \infty)$, decreciente en $(-\infty, 2)$.
  \item a) $f'(x) = 3x^2 - 3$; puntos críticos $x = \pm 1$. $x = -1$ máximo, $x = 1$ mínimo. b) Creciente en $(-\infty, -1) \cup (1, \infty)$; decreciente en $(-1, 1)$. c) $f''(x) = 6x$; cóncava abajo en $(-\infty, 0)$, arriba en $(0, \infty)$; punto de inflexión $(0, 0)$.
  \item a) $f'(x) = -\frac{1}{x^2} < 0$: decreciente en $(-\infty, 0)$ y en $(0, \infty)$. b) $f''(x) = \frac{2}{x^3}$: cóncava abajo en $(-\infty, 0)$, cóncava arriba en $(0, \infty)$. c) Asíntotas: $x = 0$ (vertical) e $y = 0$ (horizontal).
  \item $f''(x) = 12x^2 - 12 = 12(x^2 - 1) = 12(x - 1)(x + 1)$. Puntos de inflexión: $x = \pm 1$ (ambos con cambio de signo).
  \item Dominio $\mathbb{R}$. $f'(x) = e^{-x}(1 - x)$: creciente en $(-\infty, 1)$, decreciente en $(1, \infty)$; máximo en $x = 1$, $f(1) = e^{-1}$. $f''(x) = e^{-x}(x - 2)$: cóncava abajo en $(-\infty, 2)$, arriba en $(2, \infty)$; punto de inflexión $(2, 2e^{-2})$. Asíntota horizontal $y = 0$ cuando $x \to \infty$.
  \item Dominio: $x \neq \pm 1$. Intercepto $(0, 0)$. Asíntotas verticales: $x = 1$ y $x = -1$. Asíntota horizontal: $y = 0$. $f'(x) = \frac{-2x^2 - 2}{(x^2 - 1)^2} < 0$ siempre: decreciente en cada intervalo. $f''$ analiza concavidad; la información completa se usa para el esbozo.
  \item $f''(x) = 6x - 2a$. En $x = 1$: $6 - 2a = 0 \Rightarrow a = 3$ (y además cambia de signo, así que es punto de inflexión).
  \item $f'(x) = \frac{2x}{x^2 + 1}$: decreciente en $(-\infty, 0)$, creciente en $(0, \infty)$; mínimo en $(0, 0)$. $f''(x) = \frac{-2x^2 + 2}{(x^2 + 1)^2} = \frac{2(1 - x^2)}{(x^2 + 1)^2}$: cóncava arriba en $(-1, 1)$, abajo en $(-\infty, -1) \cup (1, \infty)$; puntos de inflexión en $x = \pm 1$.
\end{enumerate}


% ──────────────────────────────────────────────────────────
%  NUEVAS RESPUESTAS (15) — INSERCIÓN MANUAL
% ──────────────────────────────────────────────────────────
\subsection*{Sección 4: Rolle, TVM y Extremos (Ejercicios 36--43)}
\begin{enumerate}[resume=respuestas, leftmargin=*, label=\textbf{\arabic*.}]
  \item $f'(x) = 2x - 6 = 0 \Rightarrow x = 3$
  \item $f(0) = 0 = f(4)$; continua y derivable en $[0, 4]$; $f'(x) = 2x - 4 = 0 \Rightarrow c = 2$.
  \item $\dfrac{f(3) - f(1)}{3 - 1} = 4$; $f'(c) = 2c = 4 \Rightarrow c = 2$
  \item $f' = 3(x^2 - 1)$; crece en $(-\infty, -1) \cup (1, \infty)$
  \item $f' = 3x(x - 2)$: máximo local en $x = 0$ ($f(0) = 2$) y mínimo local en $x = 2$ ($f(2) = -2$).
  \item $f'' = 6x - 6 = 0 \Rightarrow x = 1$; punto de inflexión $(1, -1)$.
  \item $f'' = 12x(x - 2)$: cóncava hacia arriba en $(-\infty, 0) \cup (2, \infty)$ y hacia abajo en $(0, 2)$; inflexiones en $x = 0$ y $x = 2$.
  \item $\dfrac{f(4) - f(1)}{3} = \dfrac{1}{3}$; $f'(c) = \dfrac{1}{2\sqrt{c}} = \dfrac{1}{3} \Rightarrow c = \dfrac{9}{4}$
\end{enumerate}

\subsection*{Sección 5: Análisis de Gráficas y Retos (Ejercicios 44--50)}
\begin{enumerate}[resume=respuestas, leftmargin=*, label=\textbf{\arabic*.}]
  \item $f'(x) = 2x + 2 = 0 \Rightarrow x = -1$
  \item No tiene extremos locales: $f'(x) = 3x^2 \ge 0$; es estrictamente creciente en $\mathbb{R}$.
  \item Dominio $\mathbb{R}$; cortes $(1, 0)$, $(3, 0)$, $(0, 3)$; vértice $(2, -1)$ (mínimo); cóncava hacia arriba.
  \item Mínimo absoluto $-1$ (en $x = 2$); máximo absoluto $8$ (en $x = 5$).
  \item Asíntota vertical $x = 1$; asíntota oblicua $y = x + 1$; $f' = \frac{x(x - 2)}{(x - 1)^2}$: máximo local en $x = 0$, mínimo local en $x = 2$.
  \item $f' = 3(x - 1)(x - 3)$: crece en $(-\infty, 1) \cup (3, \infty)$, decrece en $(1, 3)$; máximo local $(1, 5)$, mínimo local $(3, 1)$; $f'' = 6(x - 2)$: cóncava hacia abajo en $(-\infty, 2)$, hacia arriba en $(2, \infty)$; punto de inflexión $(2, 3)$.
  \item Cuadrado de $10 \times 10$ m; área máxima $100$ m$^2$.
\end{enumerate}
\fi
\end{document}
